\chapter*{Afterword}
\phantomsection
\addcontentsline{toc}{chapter}{Afterword}
\markboth{Afterword}{}
\thispagestyle{plain}

%============================================================
\section{Reading Guide: How to Face Those "Pretentious" Terms}
%===========================================================

As a reader, you may have found yourself more than once wanting to grumble at the screen: "This author is really showing off!" That is a perfectly normal reaction---the book does occasionally contain expressions that are not exactly gentle on the brain.

\vspace{1em}

For example, consider this passage, capable of making anyone give up within three seconds:

\begin{myquote}

\textit{"All learning begins with bias. This sounds like a criticism, but it is a statement of fact. In machine learning, this is called inductive bias."}
\end{myquote}

In plain language, what it's really saying is:

\textbf{Inductive bias (learning with a certain preconception):} Any system capable of learning cannot face the world with zero assumptions. Convolutional networks assume "nearby relationships matter more"; Transformers assume "distant elements can also be related to each other."
These assumptions are not necessarily flaws---they are precisely the prerequisites that make learning possible in the first place. What this statement really means is: a model does not learn about the world from a blank slate. It always carries a certain way of seeing the world. The question is not "whether there is bias," but whether the bias is appropriate.

\vspace{1em}

Similarly, the book occasionally features intimidating-looking pseudo-formulas. If you take a closer look, all this expression is trying to say is: raw data undergoes \textbf{pruning} and \textbf{alignment}, step by step becoming \textbf{clean and aligned} data.

\begin{myquote}
    \begin{equation*}
    p_{\text{raw}}(y|x) \xrightarrow{\text{pruning}} p_{\text{clean}}(y|x) \xrightarrow{\text{alignment}} p_{\text{aligned}}(y|x)
    \end{equation*}
\end{myquote}

Why preserve this cognitive load?


While this cognitive load was not deliberately imposed, I have no intention of eliminating it entirely, for three reasons:

\textbf{Building frameworks:} The barrier to mastering individual knowledge points has dropped to near zero---any topic can be looked up with GPT at any time, provided we know what the concept is called. The role of terminology is to build the scaffolding of understanding. Details can be filled in by AI, but the big picture must live in our own minds.

\textbf{Creating an index:} Each term is a key. Once we have mastered these indices, we can precisely leverage GPT for deep learning whenever we want to explore a particular field in the future.

\textbf{Eliminating fear:} Frankly, familiarity with these terms does make us appear more knowledgeable in social settings. More importantly, once we see through the "seemingly profound" shell, we discover that the core ideas are actually not that complicated---the underlying logic of AI is not as inaccessible as it seems.

In the AI era, mastering any isolated piece of knowledge is no longer difficult. What is truly scarce is a holistic cognitive framework. The reason this book retains some professional terminology is not to \textbf{show off}, nor to \textbf{drill into} formulas, but to help readers establish a coordinate system.

\vspace{1em}

Once we build a holistic cognitive framework, the walls formed by terminology become signposts on our map.
